\pg{347} \textbf{69. O ente finito é contingente; o ente contingente é finito.}

a) \emph{O ente finito é contingente}. --- Ente \emph{finito} é aquelle que tem o ser limitado. Ora o ente, que tem o ser limitado, não é o proprio ser subsistente, não existe pela exigencia da sua essencia, é \emph{contingente}. Logo o ente \emph{finito} é \emph{contingente}.

b) \emph{O ente contingente é finito}. --- Ente \emph{contingente} é o que é por si indifferente para a existencia. Ora o ente, que é por si indifferente para a existencia, é composto de potencia e de acto, e por isso é \emph{finito}. Logo o ente \emph{contingente} é \emph{finito} \footnote{\emph{Finito} e \emph{contingente} são termos, que se correspondem. --- O \emph{finito} é composto de acto e de potencia (pois a limitação é potencia, é imperfeição), não é o proprio ser subsistente, não existe pela exigencia da sua essencia, é indifferente para a existencia e para a não-existencia, é \emph{contingente}. --- O \emph{contingente}, pela sua vez, é composto de essencia e de existencia, de potencia e de acto, e por isso é \emph{finito}; pois cada parte é finita, e, como o finito, addicionado ao finito, só dá o finito, a somma das partes é tambem finita.}.

70. A existencia dos entes finitos suppõe a \textbf{existencia} do Ente infinito.

a) \emph{Os entes finitos}, porque são \emph{entes}, possuem o ser, e, porque são \emph{finitos}, não são o proprio ser, isto é, não possuem o ser por si mesmos, pela exigencia da sua essencia. Não possuindo o ser por si mesmos, os entes finitos devem tel-o recebido de outro ente; aliás haveria um effeito sem causa proporcionada. Esse outro ente, ou é tambem finito, ou é infinito. Se é infinito, temos o que queriamos. Se é finito, subiremos \textbf{ainda, até se encontrar} necessariamente o ente, que possua o \textbf{ser} \textbf{por si mesmo}, que seja a plenitude do ser, do qual todos os \textbf{entes finitos} receberam o ser, que seja \emph{infinito}; porque, se não existisse o ente infinito, não poderia existir nenhum ente finito. Logo a existencia dos entes \emph{finitos} suppõe a existencia do Ente \emph{infinito} \footnote{O argumento é claro e convincente. Todo o ente, que existe, --- ou existe \emph{por si}, pela sua essencia, e é \emph{ens a se}. --- ou existe \emph{por outro ente}, por participação, e é \emph{ens ab alio}. Não se dá meio termo. Ora o ente \emph{finito}, como tal, não é ente \emph{por si}, isto é, não existe pela exigencia da sua essencia, mas é ente por participação, porque a sua essencia não é a sua existencia, não é o seu ser. Logo existe por outro ente, que lhe communicou a existencia. Este outro ente deve ser, em ultima analyse, o proprio ser subsistente, o ente infinito. Se não existisse o infinito, o ente por essencia, não poderia existir o finito, o ente por participação: assim como, se não existisse o calor, não poderia existir nenhum corpo quente, e, se não existisse a luz, não poderia existir nenhum corpo luminoso. Diz S. Thomaz: «Quod per essentiam dicitur est causa omnium, quae per participationem dicuntur, sicut\pgn{348}ignis est causa omnium ignitorum in quantum hujusmodi. Deus autem est ens per essentiam suam, quia est ipsum esse; omne autem aliud ens est ens per participationem, quia ens quod sit suum esse non potest esse nisi unum. Deus igitur est causa essendi omnibus aliis» (\emph{C}. \emph{Gent}., II, 15).}.
